%!TEX ROOT = thesis.tex

% =========================================================
% CHAPTER 4: LITERATURE REVIEW
% =========================================================
\chapter{LITERATURE REVIEW}
\label{chap:litreview}

This chapter reviews the scientific and technical literature that supports the design choices of this thesis. The purpose is not to claim that the proposed model replaces all existing cardiac segmentation methods. Instead, the review identifies the specific research context in which a lightweight Mamba-inspired U-Net is meaningful: automated cardiac MRI segmentation requires accurate anatomical delineation, but practical models must also consider parameter count, memory usage, and reproducible validation.

\section{Cardiac MRI Segmentation as a Clinical Computing Problem}
Cardiac Magnetic Resonance (CMR) imaging is widely used for non-invasive assessment of cardiac morphology and function. In the ACDC benchmark, the segmentation targets are the Right Ventricle (RV), Myocardium (MYO), and Left Ventricle (LV), which are clinically important because they support the computation of ventricular volumes, myocardial mass, and ejection fraction \citep{bernard2018deep}. These quantities are not obtained directly from a neural network classifier; they require spatial masks from which anatomical volumes can be computed. Therefore, segmentation quality is central to the reliability of downstream cardiac function analysis.

CMR segmentation is challenging because the difficulty is not uniform across all structures. LV segmentation is often more stable because the LV cavity is relatively compact and has a clearer blood-pool boundary. In contrast, the myocardium is a thin ring-like structure whose boundaries can be affected by partial volume effects, papillary muscles, and signal ambiguity. The RV is also difficult because it has a more variable crescent-like shape and thinner wall. These anatomical observations explain why a single aggregate Dice score is insufficient; class-specific performance should also be reported when evaluating cardiac segmentation models.

The ACDC dataset is useful for this thesis because it includes multiple cardiac conditions, including normal subjects and pathological groups such as dilated cardiomyopathy, hypertrophic cardiomyopathy, myocardial infarction, and abnormal right ventricle cases \citep{bernard2018deep}. This heterogeneity makes the benchmark more relevant than a dataset containing only healthy cardiac anatomy. Nevertheless, the ACDC training cohort remains a limited dataset, so patient-level separation between training and validation is essential. If frames from the same patient appear in both training and evaluation, the reported performance may reflect memorization of patient-specific anatomy rather than generalization.

\section{Traditional Medical Image Segmentation Methods}
Before deep learning became dominant, medical image segmentation commonly relied on intensity thresholding, deformable models, graph-based optimization, and level-set methods. These approaches remain important historically because they clarify why learning-based segmentation became attractive.

A simple thresholding method assumes that the target object can be separated from the background by intensity. Otsu's method, for example, selects a threshold $T$ by maximizing the between-class variance \citep{otsu1979threshold}. If the two classes have probabilities $\omega_0(T)$ and $\omega_1(T)$, and means $\mu_0(T)$ and $\mu_1(T)$, the objective can be written as:
\begin{equation}
    \sigma_B^2(T) = \omega_0(T)\omega_1(T)\left[\mu_0(T)-\mu_1(T)\right]^2.
\end{equation}
This formulation is elegant because it does not require annotated training data. However, it is not sufficient for cardiac MRI segmentation. MRI intensities are not standardized in the same way as CT Hounsfield Units, and myocardium, blood pool, and surrounding tissues may overlap in intensity. As a result, a global threshold cannot reliably separate RV, MYO, and LV across patients.

Active contour models, introduced by \citet{kass1988snakes}, formulate segmentation as an energy minimization problem. A contour evolves under internal smoothness forces and external image forces. A simplified energy functional can be written as:
\begin{equation}
    E_{snake} = \int_0^1 \left[\frac{1}{2}\alpha |v'(s)|^2 + \frac{1}{2}\beta |v''(s)|^2 + E_{image}(v(s))\right]ds.
\end{equation}
Here, the first two terms regulate contour continuity and smoothness, while $E_{image}$ attracts the contour to image features such as edges. This is useful when boundaries are clear, but cardiac MRI boundaries can be weak or ambiguous at basal and apical slices. Moreover, deformable models often require careful initialization, which limits their practicality for fully automated pipelines.

Level-set methods represent the contour implicitly as the zero level of a higher-dimensional function $\phi$ \citep{osher1988fronts}. A typical evolution equation is:
\begin{equation}
    \frac{\partial \phi}{\partial t} + F|\nabla \phi| = 0.
\end{equation}
The implicit representation allows topology changes more naturally than parametric snakes. However, iterative numerical optimization and sensitivity to image forces remain practical limitations. For this reason, traditional methods are useful background, but they do not remove the need for robust data-driven segmentation in heterogeneous CMR datasets \citep{chen2020cardiacreview}.

\section{Deep Learning Segmentation Architectures}
\subsection{Fully Convolutional Networks and U-Net}
Fully Convolutional Networks (FCNs) showed that classification-style neural networks could be reformulated for dense pixel-wise prediction by replacing fully connected layers with convolutional operations \citep{long2015fully}. This idea was important because segmentation requires a prediction for every pixel or voxel rather than one image-level label.

U-Net became one of the most influential architectures in biomedical image segmentation \citep{ronneberger2015u}. Its main contribution is the encoder-decoder structure with skip connections. The encoder gradually reduces spatial resolution to learn semantic features, while the decoder restores resolution to produce a segmentation mask. Skip connections transfer high-resolution spatial features from the encoder to the decoder, which helps recover fine boundaries. If $\mathbf{X}^{(l)}_{enc}$ is an encoder feature map and $\mathbf{X}^{(l+1)}_{dec}$ is a decoder feature map from a coarser level, a simplified skip-connection operation can be written as:
\begin{equation}
    \mathbf{X}^{(l)}_{\mathrm{dec}}
    =\operatorname{Conv}\!\left(
      \operatorname{Concat}\!\left[
        \mathbf{X}^{(l)}_{\mathrm{enc}},
        \operatorname{Upsample}\!\left(\mathbf{X}^{(l+1)}_{\mathrm{dec}}\right)
      \right]
    \right),
\end{equation}
where \(\operatorname{Upsample}\) restores spatial resolution,
\(\operatorname{Concat}\) joins feature channels, and
\(\operatorname{Conv}\) performs convolutional refinement. Writing the
operations by name avoids any ambiguity about the order: up-sampling occurs
before channel concatenation, and convolution follows the concatenation. This
equation explains why U-Net is suitable for medical segmentation: the decoder
does not reconstruct boundaries only from compressed semantic features; it
also receives high-resolution information from earlier encoder layers.

For volumetric data, 3D U-Net extends the same principle from 2D images to 3D tensors \citep{cicek20163d}. V-Net also used volumetric convolution and Dice-based optimization for 3D medical segmentation \citep{milletari2016v}. These works are directly relevant to this thesis because the proposed Nano-Mamba U-Net uses a 3D U-Net-style encoder-decoder as its architectural base. The contribution of this thesis is not the invention of the U-Net framework itself. Rather, the project studies whether the bottleneck of such a framework can be modified into a compact sequence-processing module.

\begin{figure}[hbt!]
\centering
\begin{tikzpicture}[
    node distance=1.25cm and 1.4cm,
    block/.style={rectangle, draw=black!80, thick, rounded corners=2pt, minimum width=3.0cm, minimum height=0.85cm, align=center},
    enc/.style={block, fill=blue!8},
    dec/.style={block, fill=green!8},
    bot/.style={block, fill=red!8},
    arrow/.style={->, >=stealth, thick},
    skip/.style={->, >=stealth, dashed, thick, draw=blue!70}
]
\node[enc] (e1) {Encoder Level 1\\High resolution};
\node[enc, below=of e1] (e2) {Encoder Level 2\\Down-sampled};
\node[bot, below=of e2] (b) {Bottleneck\\Compressed features};
\node[dec, right=3.5cm of e2] (d2) {Decoder Level 2\\Up-sampled};
\node[dec, right=3.5cm of e1] (d1) {Decoder Level 1\\Segmentation detail};
\draw[arrow] (e1) -- (e2);
\draw[arrow] (e2) -- (b);
\draw[arrow] (b) -- (d2);
\draw[arrow] (d2) -- (d1);
\draw[skip] (e2) -- node[above, font=\scriptsize] {skip} (d2);
\draw[skip] (e1) -- node[above, font=\scriptsize] {skip} (d1);
\end{tikzpicture}
\caption{Conceptual encoder-decoder structure of U-Net-style segmentation models. Encoder blocks learn increasingly abstract features, decoder blocks recover spatial resolution, and skip connections preserve boundary information.}
\label{fig:unet_concept}
\end{figure}

\subsection{Attention U-Net}
Attention U-Net introduced attention gates to suppress irrelevant encoder features before they are passed through skip connections \citep{oktay2018attention}. This is relevant in medical images because large regions of an image may contain background structures unrelated to the target anatomy. In a simplified attention gate, an encoder feature $x$ and a decoder gating signal $g$ are projected into a shared representation and transformed into an attention coefficient:
\begin{equation}
    \alpha = \operatorname{sigmoid}\!\left(
      \boldsymbol{\psi}^{\mathsf{T}}
      \delta\!\left(W_xx + W_gg + b\right)
    \right).
\end{equation}
The filtered feature can then be represented as:
\begin{equation}
    \hat{x}_i=\alpha_i x_i,
\end{equation}
where index \(i\) runs over the spatial feature elements. Thus,
\(\operatorname{sigmoid}\) maps the attention score to a coefficient,
\(\delta\) is a nonlinear activation, and \(\alpha_i x_i\) is explicit
element-wise multiplication. The key idea is not that attention automatically
guarantees better segmentation, but that it provides a mechanism for feature
selection. In the final experiment of this thesis, Attention U-Net did not
outperform the other models, which illustrates why empirical comparison under
the same split is necessary.

\subsection{Residual Networks and SegResNet}
Residual learning was introduced to make deeper networks easier to optimize \citep{he2016deep}. Instead of asking stacked layers to learn a full mapping $H(x)$, a residual block learns a residual function $F(x)$ and adds the input back to the output:
\begin{equation}
    y = F(x) + x.
\end{equation}
This identity shortcut improves gradient flow and can make deeper convolutional models more trainable. SegResNet-style medical segmentation architectures use residual blocks in an encoder-decoder setting \citep{myronenko2019segresnet}. In this thesis, the MONAI SegResNet configuration with 16 initial filters serves as a residual convolutional alternative to the proposed Nano-Mamba U-Net. The final results show that SegResNet16 achieved the highest mean Dice, which must be acknowledged rather than hidden.

\section{Transformers and Sequence Modeling in Medical Imaging}
Transformers introduced self-attention as a way to model long-range dependencies \citep{vaswani2017attention}. Vision Transformer adapted this idea to image classification by treating image patches as tokens \citep{dosovitskiy2020image}. In medical imaging, UNETR used transformer encoders for 3D segmentation \citep{hatamizadeh2022unetr}, while Swin Transformer and Swin UNETR introduced hierarchical shifted-window attention to reduce the cost of global attention \citep{liu2021swin,hatamizadeh2022swin}.

The computational concern with standard self-attention is that the attention matrix grows quadratically with the number of tokens. For an input sequence of length $N$, the attention matrix has size $N\times N$, and the simplified attention operation is:
\begin{equation}
    \operatorname{Attention}(Q,K,V)
    =\left[
      \operatorname{softmax}\!\left(
        \frac{QK^{\mathsf{T}}}{\sqrt{d_k}}
      \right)
    \right]V.
\end{equation}
Here, \(d_k\) is the key-vector dimension. The square brackets make the
scope explicit: softmax acts on the complete scaled score matrix
\(QK^{\mathsf{T}}/\sqrt{d_k}\), and the resulting attention-weight matrix is
then multiplied by \(V\). This equation shows why sequence length matters: the
multiplication \(QK^{\mathsf{T}}\) compares every token with every other token.
Window-based methods reduce this cost, but the broader issue remains important
for volumetric segmentation because 3D images can produce many spatial tokens.

It would be inaccurate to claim that all transformer models are infeasible or that they are unsuitable for medical imaging. Prior work shows that transformer-based segmentation can be effective. The more careful motivation for this thesis is narrower: for a master's project focused on lightweight deployment and reproducible experiments on a consumer GPU, it is reasonable to investigate alternatives that attempt to combine broader context modeling with lower parameter count.

\section{State Space Models and Mamba-Inspired Design}
State Space Models (SSMs) represent sequence dynamics through a hidden state. A continuous-time linear state space model can be written as:
\begin{equation}
    h'(t)=Ah(t)+Bx(t), \qquad y(t)=Ch(t).
\end{equation}
Here, $h(t)$ is the hidden state, $x(t)$ is the input signal, $y(t)$ is the output signal, and $A$, $B$, and $C$ are state transition and projection matrices. In the context of deep learning, the appeal of SSMs is that they can model long sequences without explicitly constructing a full $N\times N$ attention matrix.

Mamba extends this line of work through selective state space modeling and hardware-aware sequence processing \citep{gu2023mamba}. Recent biomedical segmentation studies have also explored Mamba-style modules for medical image segmentation. U-Mamba combines convolutional processing with state-space blocks, while SegMamba studies long-range sequential modeling for 3D medical segmentation \citep{ma2024umamba,xing2024segmamba}. VM-UNet uses Visual State Space blocks in a U-shaped architecture, Mamba-UNet uses a VMamba-based encoder-decoder and reports experiments on ACDC and Synapse, and LightM-UNet uses residual Vision Mamba layers to pursue a lightweight design \citep{ruan2024vmunet,wang2024mambaunet,liao2024lightmunet}. These works make Mamba/SSM segmentation directly relevant to this thesis, but they are related methods rather than evidence that every Mamba-inspired block will automatically improve cardiac segmentation.

However, this thesis must distinguish the original Mamba implementation, the Visual State Space or SSM blocks in those recent segmentation models, and the implemented project model. Nano-Mamba serializes compressed 3D features and applies only kernel-3 depthwise sequence convolution, token-wise scalar gating, projections, and a residual connection. It has no recurrent state-space transition, selective scan, attention matrix, or direct global all-token interaction. Its fixed raster order also creates boundary adjacencies that are not isotropic 3D neighbours. It is therefore more accurate to describe the contribution as a local Mamba-inspired gated bottleneck exploration, not as a reproduction of Mamba or the cited Visual State Space architectures.

This distinction matters for scientific integrity. If the implementation does not include the full selective scan algorithm, the thesis should not claim that it has implemented the complete Mamba architecture. The contribution is still valid, but it should be framed correctly: the project investigates whether a compact sequence-inspired bottleneck can provide a useful accuracy-efficiency trade-off in cardiac MRI segmentation.

\section{Full-Cine Temporal Segmentation and Cardiac Motion Analysis}
Cardiac cine analysis can use time in fundamentally different ways. \citet{qin2018jointmotion} jointly learned cardiac MR segmentation and motion estimation using recurrent spatial transformers, thereby producing explicit inter-frame motion information in addition to masks. \citet{yan2019cineflow} used optical-flow learning to add the temporal dimension to cine MRI analysis. Such approaches seek a deformation field or pixel/tissue correspondence between phases; when appropriately validated, this supports motion tracking beyond independent mask measurements.

The present thesis addresses a narrower downstream problem. The trained Nano-Mamba backbone receives one 3D phase at a time and has no temporal input. After all phases are inferred, a fixed circular filter combines class probabilities as $0.25P_{t-1}+0.50P_t+0.25P_{t+1}$, including wrap-around at the first and last phases. The resulting masks support LV/RV/MYO volume curves, phase timing, EF, LV-centroid displacement, and a myocardial radial-distance summary. These are segmentation-derived global temporal and motion surrogates. They do not establish point correspondence, a dense deformation field, optical flow, local myocardial displacement, or strain. This distinction defines both the contribution and the validation boundary of the full-cine stage.

\section{Loss Functions and Evaluation Metrics}
Medical segmentation is often class-imbalanced because background voxels dominate the image volume, while structures such as myocardium occupy a smaller region. Dice-based objectives are commonly used because they directly optimize overlap between prediction and ground truth \citep{milletari2016v,sudre2017generalised}. For a predicted binary mask $P$ and ground-truth mask $G$, the Dice Similarity Coefficient is:
\begin{equation}
    \mathrm{DSC}(P,G)=\frac{2|P\cap G|}{|P|+|G|}.
\end{equation}
In multi-class segmentation, Dice can be computed separately for RV, MYO, and LV, and then averaged. This is the reporting strategy used in this thesis.

Cross-entropy is also widely used because it provides voxel-wise classification supervision. If $y_i$ is the true class at voxel $i$ and $p_{i,c}$ is the predicted probability for class $c$, the multi-class cross-entropy can be expressed as:
\begin{equation}
    L_{\mathrm{CE}}=-\sum_i\sum_c \mathbf{1}(y_i=c)\log p_{i,c}.
\end{equation}
The final experiment uses MONAI's Dice-Cross Entropy loss, combining region-overlap supervision with voxel-wise classification supervision. This choice is practical rather than novel; it follows established segmentation practice and is supported by the MONAI framework \citep{cardoso2022monai}.

Evaluation should also avoid leakage. If the same patient contributes samples to both training and validation, the measured Dice may overestimate generalization. Therefore, the final experimental design of this thesis uses patient-level splitting. This methodological point is as important as the architecture itself, because a strong model comparison is only meaningful when all models are trained and evaluated under the same split.

\section{Summary of the Literature Gap}
The literature shows that U-Net-style models remain strong foundations for biomedical segmentation, attention mechanisms can improve feature selection, residual models provide strong convolutional baselines, and transformer/SSM approaches offer explicit mechanisms for longer-range modeling. Recent Mamba segmentation studies include full Visual State Space or state-space components, whereas the present block is a substantially simpler local gated operation. The first gap addressed here is therefore not the absence of segmentation models, but the need to evaluate this compact implementation honestly under a reproducible patient-level protocol.

The motion literature also distinguishes segmentation-derived trajectories from learned motion fields and optical flow. The second gap addressed here is a bounded, auditable bridge from one validated spatial checkpoint to complete-cine global function and motion summaries, together with a direct comparison of frame-wise inference and fixed circular probability fusion. Based on this review, Nano-Mamba U-Net is positioned as a lightweight Mamba-inspired extension of a 3D U-Net-style model embedded in a complete but non-dense temporal analysis. The appropriate questions are whether it offers a useful accuracy--efficiency trade-off on ACDC and what its segmentation outputs can support without claiming learned temporal representation or tissue tracking.
